\chapter{{\fontsize{14pt}{21pt}\selectfont\MakeUppercase{Приложение В. Руководство пользователя мобильного приложения и руководство администратора веб-панели}}}
\label{app:user-admin}

\section*{В.1 Назначение}

Настоящее приложение описывает пошаговые действия слушателя в мобильном клиенте и администратора или преподавателя в веб-панели при работе со стендом, собранным по инструкции приложения~Г\cite{flutter_docs,react_docs,fastapi_docs}.

Общая политика демонстрации. Пароли в демонстрационных аккаунтах выражены литерально исключительно в учебных целях; в эксплуатации заменять на управляемые секретные фразы.

\section*{В.2 Демонстрационные учётные записи}

\begin{table}[H]
\centering
\caption{Демонстрационные учётные записи программного комплекса}
\label{tab:appv-demo-users}
{\setstretch{1}\fontsize{12pt}{14pt}\selectfont
\begin{tabular}{|l|l|p{72mm}|}
\hline
Электронная почта & Пароль (демо) & Назначение роли в сценариях поверки \\
\hline
sysadmin@example.com & Password123! & высокоуровневое администрирование платформы \\
\hline
admin@acme.example.com & Password123! & администрирование организации-арендатора учебного кейса \\
\hline
teacher@acme.example.com & Password123! & подготовка курсов, тестирования и аналитический просмотр \\
\hline
learner1@acme.example.com & Password123! & прохождение уроков, адаптивного теста и получение рекомендаций \\
\hline
\end{tabular}
}
\end{table}

\section*{В.3 Руководство пользователя мобильного приложения}

\subsection*{В.3.1 Подключение к серверу}

\begin{enumerate}
  \item Открыть проект каталогом mobile, выполнить flutter pub get.
  \item При использовании эмулятора Android зафиксировать базовый URL API как http://10.0.2.2:8000/api/v1; при физическом устройстве --- IP узла локальной сети с портом 8000.
  \item Запуск: выполнить flutter run {-}{-}dart-define=API\_BASE=http://10.0.2.2:8000/api/v1 на эмуляторе Android; на устройстве подставить IP узла сети вместо 10.0.2.2.
\end{enumerate}

\subsection*{В.3.2 Вход слушателя}

\begin{enumerate}
  \item Открыть экран входа, ввести learner1@acme.example.com и Password123!.
  \item После успешной аутентификации убедиться, что назначенные курсы отображаются в списках и доступны для открытия.
\end{enumerate}

\subsection*{В.3.3 Обучение и адаптивный тест}

\begin{enumerate}
  \item Из списка курсов открыть назначенный курс, перейти к урокам, пролистать страницы содержания; при наличии видеоматериала задействовать стандартные элементы воспроизведения плеера.
  \item Для контроля знаний открыть сценарий адаптивного теста (экран TestFlowScreen); подтвердить старт попытки, поочерёдно отвечать на вопросы, наблюдая возможное изменение уровня сложности между вопросами.
  \item Завершить попытку, просмотреть текстовые рекомендации по темам после обработки ответов системой.
\end{enumerate}

\section*{В.4 Руководство администратора организации или преподавателя (веб-панель)}

\subsection*{В.4.1 Подключение к веб-службе}

\begin{enumerate}
  \item По умолчанию при контейнеризированном стенде панель открывается на http://localhost:8081.
  \item Убедиться, что при сборке веб-панели задаётся корректный базовый URL серверной части с суффиксом /api/v1.
\end{enumerate}

\subsection*{В.4.2 Проверочный сценарий администратора организации}

\begin{enumerate}
  \item Авторизация аккаунтом admin@acme.example.com.
  \item Проверка раздела пользователей: назначить роли, при необходимости создать учебную группу для массового назначения.
  \item Создание объекта учебной программы или проверка существования демонстрационной программы после сидирования.
\end{enumerate}

\subsection*{В.4.3 Проверочный сценарий преподавателя}

\begin{enumerate}
  \item Авторизация аккаунтом teacher@acme.example.com.
  \item В разделе «Тесты» подготовить тестируемый объект или отредактировать существующий: поля ограничителя попытки, активный тест, банк вопросов.
  \item Назначить курс аудитории слушателей через соответствующий раздел.
  \item Открыть аналитический раздел панели (сводка по активности обучения) и проверить отображение агрегированных результатов после прохождения тестирования слушателем.
\end{enumerate}

\subsection*{В.4.4 Системная роль платформы (по необходимости)}

Аккаунт sysadmin@example.com служит контрольным наблюдением за экземпляром при многофирменной модели данных; использовать только на защите по сценарию, согласованному с методическим указанием.

\section*{В.5 Известные ограничения демонстрации}

При работе стенда только в локальной сети необходимо, чтобы браузер и мобильный клиент разрешали обращение по протоколу HTTP к адресам узла сборки; в производственном контуре использовать обратное проксирование с TLS.
